\section{Prospective Evaluation Protocol}
\label{sec:evaluation}

A useful evaluation must distinguish three propositions: the Episode
representation preserves information unavailable in flattened histories; the
discovery mechanism uses that information better than conservative baselines;
and the resulting candidate transfers without exporting more complexity or
governance harm than it compresses.

\subsection{Protocol A: representation ablation}

Construct several immutable corpora from the same completed work Episodes:

\begin{enumerate}[leftmargin=*]
  \item full Episode assets with perspective, causality, responsibility,
        evidence, gaps, and stewardship metadata;
  \item provenance-preserving events without participant consequence and
        ontology declarations;
  \item anonymous event sequences with timestamps and payloads; and
  \item endpoint Fact cuts only.
\end{enumerate}

Blind evaluators to the known candidate names. Ask them to reconstruct why
specified decisions differed, identify repeated unnamed burdens, and propose
or reject candidate distinctions. Measure reconstruction time, unresolved
questions, evidence alignment, candidate yield, false-candidate rate, and
inter-evaluator agreement. If the full Episode condition does not improve
recoverability or reduces it by imposing excessive structure, RQ1 is weakened.

The ablation must hold source cases and resource budgets constant. Otherwise a
larger corpus or more capable model could be mistaken for a representation
effect.

\subsection{Protocol B: preregistered shadow KFD-6 loop}

Before new work begins, freeze $E_n$ and $O_n$, preregister method procedures,
budgets, metrics, falsifiers, and autonomy boundaries, and reserve a held-out
cut. Run at least two method families, for example perspective transformation
and anomaly/reconstruction analysis, plus fixed-ontology and no-new-Primitive
baselines.

Candidate names and attractive narratives are not the primary endpoint. The
endpoint is whether a candidate survives KFD-5 minimum closure, alternative,
deletion, fuse, falsifier, and dogfood gates and then improves a held-out
decision. Retain rejected candidates and score why they failed. A system that
generates many plausible objects but cannot reject them fails the experiment.

Shadow mode prohibits the loop from changing production ontology or participant
responsibility. An independent evaluator receives the held-out evidence. A
separate accountable authority decides whether any result warrants a bounded
intervention.

\subsection{Protocol C: transfer and intervention}

A provisional candidate that passes Protocol B is mapped into at least one
context outside its founding workstream. The test permits different names and
representations. It asks whether an equivalent distinction:

\begin{itemize}[leftmargin=*]
  \item reduces repeated reconstruction or coordination cost;
  \item changes a decision that remained ambiguous under the prior ontology;
  \item preserves source Fact, Episode, perspective, and authority boundaries;
  \item admits a smaller rival representation if one works better; and
  \item remains removable or reversible when consequences disagree with the
        qualification.
\end{itemize}

Transfer within one organization's repositories is weaker than independent
adoption. Both are useful but must not share a label.

\subsection{Protocol D: governance and stewardship}

The experiment audits the corpus itself. For sampled Episodes, verify integrity
roots, retention basis, consent or other legitimate capture boundary,
redaction, deletion behavior, export/import fidelity, provider migration, and
the distinction between authoritative records and derived projections. Invite
consequence-bearing participants to challenge replays attributed to them.

Governance failure can invalidate a technically promising result. A candidate
is not qualified if producing it requires undeclared surveillance, prevents
applicable deletion, silently broadens purpose, or makes independent challenge
impossible.

\subsection{Decision rule}

The strong hypothesis survives only if the Episode condition produces a
candidate that passes KFD-5, improves at least one held-out decision over both
conservative baselines, rejects attractive false candidates, transfers across
a declared boundary, and preserves governance constraints at acceptable cost.
Failure should narrow the representation, mechanism, or domain claim. It must
not be repaired by relabeling retrieval or clustering as Primitive discovery.
